\chapter*{Список сокращений и обозначений}
\addcontentsline{toc}{chapter}{Список сокращений и обозначений}
\label{ch:abbreviations}

\begin{description}[
    style=nextline,
    leftmargin=2.2cm,
    labelwidth=1.6cm,
    labelsep=0.4cm,
    font=\normalfont\bfseries,
    itemsep=0.15cm
]
    \item[BC]   Behavioral Cloning (обучение методом поведенческого клонирования)
    \item[CLIP] Contrastive Language--Image Pre-training
    \item[DoF]  Degrees of Freedom (степени свободы)
    \item[DPO]  Direct Preference Optimization
    \item[FAST] Frequency-space Action Sequence Tokenization (частотная токенизация действий)
    \item[GRPO] Group Relative Policy Optimization
    \item[LLM]  Large Language Model (большая языковая модель)
    \item[LoRA] Low-Rank Adaptation (низкоранговая адаптация)
    \item[LOOP] Leave-One-Out PPO (безкритиковая схема RL)
    \item[MDP]  Markov Decision Process (марковский процесс принятия решений)
    \item[MLP]  Multilayer Perceptron (многослойный перцептрон)
    \item[OFT]  Optimized Fine-Tuning (оптимизированный рецепт дообучения OpenVLA)
    \item[OXE]  Open X-Embodiment (набор робототехнических данных)
    \item[PPO]  Proximal Policy Optimization
    \item[RL]   Reinforcement Learning (обучение с подкреплением)
    \item[RLOO] REINFORCE Leave-One-Out
    \item[SAC]  Soft Actor-Critic
    \item[SFT]  Supervised Fine-Tuning (дообучение с учителем)
    \item[SOTA] State of the Art (наилучший достигнутый результат)
    \item[SR]   Success Rate (доля успешных выполнений)
    \item[VLA]  Vision-Language-Action (визуально-языковая модель действий)
    \item[VLM]  Vision-Language Model (визуально-языковая модель)
    \item[ДИ]   Доверительный интервал
\end{description}

\endinput
